\section{Orientations}
The idea of Morse homology is now to find an invariant under manipulations of the Morse--Smale function. Obviously neither the number of critical points nor the number of flow lines stays invariant. However certain linear combinations stay invariant, when one establishes a sophisticated way to count flow lines. For this count we want to introduce signs to the flow lines between critical points with index one apart where the signs rely on orientations of over the flow lines.
\subsection{Oriented Vector Spaces and Bundles}




\begin{definition}[Orientation]
	Let $V$ be a finite dimensional $\R$ vector space and $\det(V)\coloneq \Lambda^{\dim(V)}$ the maximal exterior algebra (or Grassmann algebra). By the principle of permanence we define $\Lambda^0(V)=\R$. An orientation of $V$ is a choice of one of the two connected components of $\det(V)\setminus \{0\}$.
\end{definition}
\begin{definition}[Convention]
    If we use $\R^n$ as an oriented vector space we fix $e\coloneq e_1\wedge \cdots \wedge e_n$ to induce the orientation.
\end{definition}
\begin{remark}
	Classically, an orientation of a vector space $V$ of dimension $m$ is given by an ordered basis $\mathfrak{b}=(b_1,\dots b_m)$ of $V$ and hence the product $b_1\wedge \dots \wedge b_m$ is an element of the maximal exterior product $\det(V)=\Lambda^mV$. Two ordered bases define different elements in $\det(V)$ and their difference is given by the determinant of the base change matrix. (This is just a calculation using multilinearity and the alternating property). Hence a choice of equivalence class of bases modulo base change matrix of positive determinant is the same as a choice of one connected component of $\det(V)\setminus \{0\}$.
\end{remark}
\begin{definition}[Determinant of an endomorphism]
	Let $f:V\to V$ be an endomorphism. Then we define 
	\begin{align}
		\Lambda^{\dim(V)}(f):\det(V)&\to \det(V)\\
v_1\wedge \dots \wedge v_{\dim(V)}&\mapsto f(v_1)\wedge \dots \wedge f(v_{\dim(V)})\,.
	\end{align} Since $\dim(\det(V))=1$, the function $\Lambda^{\dim(V)}(f)$ is given by a scaler multiplication with a nonzero element $\lambda \in \R$. We define the determinant $\det(f)\coloneq \lambda$  to be said scalar. Obviously, if $\det(V)$ is positive, then $f$ is orientation preserving in the sence, that $f$ applied to the ordered basis gives again an ordered basis that is compatible.
\end{definition}
\begin{remark}\label{rem:orient-calculateDeterminatViaMatrixDeterminant}
	From linear algebra we know, that we can calculate $\det(f)=\det(M^{\mathfrak{b}}_\mathfrak{b}(f))$, where $\mathfrak{b}$ is an ordered basis of $V$. Since the determinant of a matrix commutes with taking products, we have that the determinant is invariant under conjugation of linear endomorphisms.
\end{remark}
\begin{definition}[Orientation and Exact Sequences]\todo{FIX}
A short exact sequence of vector spaces 
\begin{center}

	\begin{tikzcd}
		0 \arrow[r] & A \arrow[r, "f"] & B \arrow[r, "g"] & C \arrow[r] & 0
	\end{tikzcd}
\end{center}
induces a canonical isomorphism 
	\begin{align*}
		 \det(A)\otimes \det(C)& \to \det(B)\\
		 (a_1\wedge \dots\wedge a_l) \otimes (c_1\wedge \dots \wedge c_k) &\mapsto (f(a_1)\wedge \dots f(a_l)\wedge \tilde{c_1}\wedge \dots \wedge \tilde{c_k})\, ,
	\end{align*} where $\tilde{c_j}\in g^{-1}(c_j)$ . The well-definedness of this map (invariance under chosen lift of $c$) can be calculated, since two different choices differ in an element in the kernel of $g$ which equals the image of $f$ and hence the different choices vanish in the wedge product. By this isomorphism we get an orientation induced in the middle from the outside. Notice the convention, that the basis from $A$ is placed in front of the basis from $C$. 
	Notice how the definition is consistent with the edge cases of trivial vector spaces $A$ or $C$. If $A=0$ we get an isomorphism $B\cong C$ and a basis from $C$ induced to $B$, that is either the induced one from $g$, or exactly the opposite one, if the trivial vector space $A$ is negatively oriented.
\end{definition}
\begin{remark}
	Let 
	\begin{center}
		\begin{tikzcd}
		0 \arrow[r] & A \arrow[r, "f"] & B \arrow[r, "g"] & C \arrow[r] & 0
	\end{tikzcd}
	\end{center}

be a short exact sequence of vector spaces and $\mathfrak{a}$ be a basis of $A$, $\mathfrak{c}$ be a basis of $C$ and $\tilde{\mathfrak{c}}$ be a lift of $\mathfrak{c}$ via $g$. Replacing the basis $\mathfrak{a}$ with $P\mathfrak{a}$ and $\mathfrak{c}$ with $Q\mathfrak{c}$ where $P,Q$ are invertible matrices, gives us a new induced basis on $B$, which differs from the original one by a matrix $\begin{pmatrix}
	P&*\\
	0&Q
\end{pmatrix}$, where $*$ depends on the chosen lift. Hence the corresponding elements in $\det(B)$ differ by the product $\det(P)\det(Q)$. 
\end{remark}
\begin{definition}
	For a rank $r$ vector bundle $E\to X$ we set $\det(E)\coloneq \Lambda^r(E)$. An orientation is an equivalence class of a non-vanishing section into the determinant bundle, where two sections are equivalent if they differ by multiplication of a positive scalar valued function. $E$ is orientable if such a section exists, i.e. if $\det(E)$ is trivial. An orientation of a bundle restricts to an orientation in each fibre.
\end{definition}

\begin{lemma}\label{lem:morse-OrientationsInAPointInduceOrientationsOfBundle}
	Let $E\to X$ be an orientable vector bundle and $X$ be connected. For each $x\in X$ we get the bijection induced from the restriction to the fibre:
	\begin{equation*}
		\{\text{Orientation on }E\} \to \{\text{Orientation on } E_x\}\,.
	\end{equation*} Hence, an orientation in one fibre can introduce an orientation in every other fibre.
\end{lemma}
\begin{proof}
	Two orientations over $E$ are given as nowhere-vanishing sections $s,s'$. The quotient $x\mapsto \frac{s}{s'}$ is continuous and its image is contained in $\R\setminus \{0\}$. Hence the sign of the quotient is a continuous function into $\{-1,1\}$. Since $X$ is connected we conclude that the sign of said quotient is a constant function. Thereby, if the sections induce the same orientation, they differ by multiplication of the positive number $\frac{s(x)}{s'(x)}$ and then the sign of $\frac{s}{s'}$ is globally one. The surjectivity is clear by mapping the sections $s$ and $-s$.
\end{proof}
\begin{proposition}
	Let 
	\begin{center}
    \begin{tikzcd}
        0\arrow[r] & E'\arrow[r,"f"] & E \arrow[r,"g"] & E''\arrow[r] & 0
    \end{tikzcd}
	\end{center} be a short exact sequence of orientable vector bundles. Then the fibrewise
    isomorphisms $\det(E'_p)\otimes\det(E''_p)\to\det(E_p)$ from the
    vector space case assemble into an isomorphism of line bundles
    \[
        \det(E')\otimes\det(E'')\;\xrightarrow{\;\cong\;}\;\det(E)\,.
    \]
    In particular, orientations of $E'$ and $E''$ induce a canonical
    orientation of $E$.
\end{proposition}
\begin{proof}
	Fibrewise the map is the canonical isomorphism from the vector
    space case: a lift $\tilde c\in g_p^{-1}(c)$ exists since $g_p$ is
    surjective, and the resulting element
    $f(a_1)\wedge\dots\wedge f(a_l)\wedge\tilde c_1\wedge\dots\wedge
    \tilde c_k$ is independent of the choice of lift.

	 It remains to
    check that this bundle map is continuous (resp.\ smooth) to ensure that it maps sections smoothly into sections, which is
    a local statement: locally choose a frame $(a_1,\dots,a_l)$ of $E'$
    and a frame $(c_1,\dots,c_k)$ of $E''$; since $g$ is an epimorphism
    it admits local sections, so we may choose local lifts
    $\tilde c_j$ of the $c_j$ which are (smooth) local sections of $E$.
    Then the map is given in the induced local frames of the
    determinant bundles by wedging these sections, hence continuous
    (resp.\ smooth). By this, two non-vanishing sections get mapped to a non-vanishing section and hence, orientations on $E'$ and $E''$ which are represented by sections give a orientation over $E$. 
\end{proof}
\begin{definition}
	We call a manifold orientable, if its tangent bundle is \textbf{orientable}. We call a touple $(M,[\omega])$ where $[omega]$ denotes an orientation of $TM$ an \textbf{oriented manifold}. We will often omitt the orientation and just call $M$ oriented.
\end{definition}

\begin{definition}[Orientation Preserving Charts and Maps]
Let $(M,[\omega])$ be an oriented $n$-dimensional manifold and let $\phi:U\to \R^n$ be a chart. We call $\phi$ \textbf{orientation preserving} if for every $x\in U$ $\partial_1(x)\wedge\cdots \wedge \partial_n(x)$ and $\omega(x)$ lie in the same connected component of $\det(T_xM)\setminus\{0\}$. More generalyly, let $(N,[\eta])$ be another oriented $n$-dimensional manifold and let $h:M\to N$ be differentiable at $x$ with $\diff h_x$ invertible. We define 
\begin{equation*}
    \sign \det(\diff h _x)\coloneq \begin{cases*}
        +1& if $\Lambda^n(\diff h_x)(\omega (x))$ and $\eta (x)$ induce the same orientation in $T_h(x)N$\\
        -1& otherwise.
    \end{cases*}
\end{equation*}
\end{definition}
\begin{remark}
    One easyly checks that $  \sign \det(\diff h _x)$ is well-defined and reversing both orientations doesn't affect the sign. In particular, for $N=M$ the number $  \sign \det(\diff h _x)$ does not depend on the choosen orientation. 
\end{remark}
\begin{lemma}[Computing the Sign in Charts]\label{lem:orient-computingSignInCharts}
    Let $(M,[\omega])$ and $(N,[\eta])$ be orinted $n$-dimensional manifolds and $\phi$ be an orinetation preserving chart arround $x$ in $M$ and $\psi$ one arround $h(x)$. Then
    \begin{equation*}
        \sign\det(\diff h_x)=\sign \det \left(\restr{\diff\left(\psi\circ h\circ \phi^{-1}\right)}{\phi(x)}\right)\,,
    \end{equation*} wherre the right hand side calculates the sign of the determinant of a matrix.
\end{lemma}
\begin{proof}
    Define $A\coloneq \restr{\diff\left(\psi\circ h\circ \phi^{-1}\right)}{\phi(x)}=\restr{\diff \psi}{h(x)}\circ \restr{\diff h}{x}\circ (\restr{\diff \phi}{x})^{-1}$ and denote the standart orientation of $\R^n$ by $e$. Since $\phi$ is orientation preserving and $\Lambda^n(\restr{\diff \phi}{x})(\partial_1(x)\wedge\cdots\wedge\partial_n(x))=e$, we have by linearity $\Lambda^n(\restr{\diff \phi}{x})(\omega(x))=a\cdot e$ where $a>0$. Similar $\Lambda^n(\restr{\diff \psi}{h(x)})(\eta(h(x)))=b\cdot e$ with $b>0$. Hence
    \begin{equation*}
        \Lambda^n\left(\restr{\diff \psi}{h(x)}\right)\circ \Lambda^n\left(\restr{\diff h}{x}\right)(\omega(x))
        =\Lambda^n(A)\circ \underbrace{\Lambda^n\left(\restr{\diff \phi}{x}\right)(\omega(x))}_{=a\cdot e}=a \det(A)\cdot e
    \end{equation*} Now by definition we need to check if $\Lambda^n(\diff h_x)(\omega (x))=c\cdot \eta (x)$ for $c>0$, which is equivalent to 
    \begin{equation*}
       \underbrace{ \Lambda^n\left(\restr{\diff \psi}{h(x)}\right)\circ \Lambda^n\left(\restr{\diff h}{x}\right)(\omega(x))}_{=a \det(A)\cdot e}=c\cdot \underbrace{\Lambda^n\left(\restr{\diff \psi}{h(x)}\right)(\eta (x))}_{=b\cdot e}\,.
    \end{equation*} which is given if and only if $\det(A)>0$.
\end{proof}
\begin{corollary}\label{cor:orient-existsOrientationPreservingAtlas}
    A smooth manifold $M$ of dimension $n$ is orientable, if and only it has an atlas such that the determinants of the Jacobians of the transition functions are positive.
\end{corollary}
\begin{proof}
   Let $[\omega]$ be an orientation on $M$ and let $U_\alpha$ be an open cover of $M$ with all $U_\alpha$ connected and the charts $\phi_\alpha:U_\alpha\to \R^n$. Let  $\partial_1^\alpha\wedge \cdots \wedge\partial_n^\alpha$ be the induced local section into the determinant bundle. Then 
   \begin{equation*}
   \restr{\omega}{U_\alpha}=g_\alpha \partial_1^\alpha\wedge \cdots \wedge\partial_n^\alpha
   \end{equation*} whith $g_\alpha$ non vanishing and $U_\alpha$ connected. If $g_\alpha$ is negative, replace $\phi_\alpha$ with $\tilde{\phi_\alpha}(x)=(-(\phi_\alpha(x))_1,\dots,\phi_\alpha(x)_n)$, which changes the sign og $g_\alpha$. Doing this in all charts gives a orientaion preserving atlas. Hence, by defininition $\sign \det(\restr{\diff \id_M}{x})=1$ and by Lemma \ref{lem:orient-computingSignInCharts} we can calculate $\sign \det(\restr{\diff \id_M}{x})$ in this atlass by the determinants of the Jacobians of the transition functions. Hence those determinants are positive.

 Fot the other direction we want to glue together the local sections $\partial_1^\alpha\wedge \cdots \wedge\partial_n^\alpha$ using a smooth partition of unity subordinate to the open cover $\{U_i\}$. This yealds something non-vanishing, as all local section are pointwise positive multiples of each other where they intersect. 
\end{proof}
\begin{corollary} \label{cor:orient-transitionFunctionPositiveDeterminantAlsoOrientationPreserving}
    The proof of Corollary \ref{cor:orient-existsOrientationPreservingAtlas} also shows that if one chart is orientation preserving, then also another chart that intersects with the first and has positive determinant of the Javobian of the transition function.
\end{corollary}

\begin{lemma}
    Let $M$ be a smooth manifold that can be covered with two charts such that their intersection is connected. Then $M$ is orientable.
\end{lemma}
\begin{proof}
    Just ensure that the two charts are compatible such that they build an atlas such that the determinants of the Jacobians of the transition functions are positive. To achive this one just manipulates one of the charts.
\end{proof}
\begin{corollary}
    $S^n$ for $n\geq 2$ is orientable, as it is coverd by the two stereographic projections. $S^1$ and $S^0$ are also orientable which can be seen by picking charts such that they fit into corollary \ref{cor:orient-existsOrientationPreservingAtlas}.
\end{corollary}







\subsection{K-theoretic Orientation}
K theoretic orientation relies on bases induced generators of the K group of the sphere. Hence, we want to fix a canonical map from $R^n$ to $S^n$:

\begin{definition}[From Real Space to the Sphere]\label{def:orient-fromRelaSpaceToTheSphere}
    We now want to construct a canonical homeomorphism between $S^m$ and $\overline{\R^m}$. First we use the stereographic projection from the South Pole:
    \begin{align*}
        \pi_S:S^m\setminus \{S\}\subseteq \R^{m+1} &\to       \R^m\\
        (x_1,\dots,x_m,x_{m+1})                            &\mapsto   \left(\frac{x_1}{1+x_{m+1}},\dots, \frac{x_m}{1+x_{m+1}} \right)
    \end{align*} with its inverse: 
    \begin{align*}
        \pi_S^{-1}:\R^m       &\to     S^m\setminus\{S\}\subseteq \R^{m+1}\\
        (y_1,\dots y_m)       &\mapsto  \left(\frac{2y_1}{1+\norm{y}^2},\dots,\frac{2y_m}{1+\norm{y}^2},\frac{1-\norm{y}}{1+\norm{y}} \right)
    \end{align*}
    Those are inverse maps and they give rise to maps between their one-point-compactifications. Hence $\overline{\R^m}\cong S^m$. Lets calculate $\pi_S^{-1}\circ \pi_s$, since we will later need a little bit of insight in those calculations:
    \begin{align*}
        \pi_S^{-1}\circ \pi_s(x_1,\dots ,x_{m+1})
        &= \pi_S^{-1}\left(\frac{x_1}{1+x_{m+1}},\dots, \frac{x_m}{1+x_{m+1}} \right)\\
        &= \left(\frac{2\frac{x_1}{1+x_{m+1}}}{1+\sum_{i=1}^m(\frac{x_i^2}{(1+x_{m+1})^2})},\dots, \frac{2\frac{x_m}{1+x_{m+1}}}{1+\sum_{i=1}^m(\frac{x_i^2}{(1+x_{m+1})^2})},\frac{1-\sum_{i=1}^m(\frac{x_i^2}{(1+x_{m+1})^2})}{1+\sum_{i=1}^m(\frac{x_i^2}{(1+x_{m+1})^2})} \right)\\
    \end{align*} Before we keep calculating, we notice the denominator can be rewritten as:
    \begin{align*}
        1+\sum_{i=1}^m(\frac{x_i^2}{(1+x_{m+1})^2})=\frac{(1+x_{m+1})^2+\sum_{i=1}^m x_i^2}{(1+x_{m+1})^2}
    \end{align*} And hence:
    \begin{align}
        \frac{2\frac{x_j}{1+x_{m+1}}}{1+\sum_{i=1}^m(\frac{x_i^2}{(1+x_{m+1})^2})}& = \frac{2\frac{x_j}{1+x_{m+1}}}{\frac{(1+x_{m+1})^2+\sum_{i=1}^m x_i^2}{(1+x_{m+1})^2}} \\
        & =\frac{2x_j(1+x_{m+1})}{(1+x_{m+1})^2+\sum_{i=1}^m x_i^2} \\
        & =\frac{2x_j(1+x_{m+1})}{(1+x_{m+1})^2+\sum_{i=1}^m x_i^2}\\
        & = \frac{2x_j(1+x_{m+1})}{1+2 x_{m+1}+\underbrace{\sum_{i=1}^{m+1} x_i^2}_{=1}} \\
        & = x_j
    \end{align}
    And the calculation of the last term is:
    \begin{align*}
        \frac{1-\sum_{i=1}^m(\frac{x_i^2}{(1+x_{m+1})^2})}{1+\sum_{i=1}^m(\frac{x_i^2}{(1+x_{m+1})^2})} 
        & = \frac{\frac{(1+x_{m+1})^2-\sum_{i=1}^m x_i^2}{(1+x_{m+1})^2}}{\frac{(1+x_{m+1})^2+\sum_{i=1}^m x_i^2}{(1+x_{m+1})^2}}\\
        & =\frac{(1+x_{m+1})^2-\sum_{i=1}^m x_i^2}{(1+x_{m+1})^2+\sum_{i=1}^m x_i^2} \\
        & = \frac{1+2x_{m+1}+x_{m+1}^2-\sum_{i=1}^m x_i^2}{1+2x_{m+1}+\underbrace{x_{m+1}^2+\sum_{i=1}^m x_i^2}_{=1}} \\
        & = \frac{1+2x_{m+1}+2x_{m+1}^2-\sum_{i=1}^{m+1} x_i^2}{2+2x_{m+1}} \\
        & = \frac{x_{m+1}(2+2x_{m+1})}{2+2x_{m+1}}=x_{m+1}
    \end{align*}
This concludes the calculation of $\pi_S^{-1}\circ \pi_s=\id$ and $\pi_s\circ\pi_S^{-1}=\id$ can be calculated in a similar manner.
\end{definition}

\begin{lemma}[Open subspaces of compact spaces]\label{lem:orient-quotientIsCompactification}
	Let $Z$ be a compact Hausdorff space and $U\subseteq Z$ open. Then the map
	\begin{equation*}
		Z\big/ (Z\setminus U)\to \overline{U}\,,\quad [z]\mapsto z \text{ for } z\in U\,,\quad [Z\setminus U]\mapsto \infty
	\end{equation*}
	is a homeomorphism.
\end{lemma}
\begin{proof}
    The map is the quotient of the function $\phi:Z\to \overline{U}$ that is the identity on $U$ and sends $Z\setminus U$ constant to $\infty$. Hence, we only need to prove continuity for $\phi$. Nor $O\subseteq \overline{U}$ is open, if it is open in $U$, or the complement of a compact set $C\subset U$ in $\overline{U}$. In both cases, $\phi^{-1}(U)$ is open and hence $\phi$ is continuous. Since the quotiont map is now a continuouos bijection from a compact space to a Hausdorff space, it is a homeomorphism.
\end{proof}

\begin{lemma}[Compactification of a product]\label{lem:orient-compactificationOfProduct}
	Let $X$ and $Y$ be locally compact Hausdorff spaces, and point $\overline{X}$ and
	$\overline{Y}$ at their respective points at infinity. Then there is a
	homeomorphism
	\begin{equation*}
		\Xi_{X,Y}:\overline{X\times Y}\xrightarrow{\ \cong\ } \overline{X}\wedge \overline{Y}
	\end{equation*}
	which is the identity on $X\times Y$ and sends $\infty$ to the base point.
	It is natural in the sense that for proper maps $f:X\to X'$ and $g:Y\to Y'$ the
	square
	\begin{center}
		\begin{tikzcd}
			\overline{X\times Y} \arrow[r, "\Xi_{X,Y}"] \arrow[d, "\overline{f\times g}"'] & \overline{X}\wedge\overline{Y} \arrow[d, "\overline{f}\wedge\overline{g}"] \\
			\overline{X'\times Y'} \arrow[r, "\Xi_{X',Y'}"']                               & \overline{X'}\wedge\overline{Y'}
		\end{tikzcd}
	\end{center}
	commutes.
\end{lemma}
\begin{proof}
    Let $Z\coloneq \overline{X}\times \overline{Y}$. Then $Z$ is a compact Hausdorfspace that contains $U\coloneq X\times Y$ as an open subset. The complement of $U$ consists of those points where one of the coordinates is infinit and hence
    \begin{equation*}
        Z\setminus U\,=\,(\{\infty\}\times \overline{Y})\cup (\overline{X}\times \{\infty\})\,=\, \overline{X}\vee \overline{Y}\,.
    \end{equation*}Therefore, by definition of the smash product, $Z\big/(S\setminus U)=\overline{X}\wedge \overline{Y}$ and by Lemma \ref{lem:orient-quotientIsCompactification} we can conculde that 
    \begin{equation*}
        \overline{X\times Y} \cong (Z\big/ (Z\setminus U)=\overline{X}\wedge \overline{Y}\,.
    \end{equation*}
    For the naturality notice that the product of proper maps is propper and hence all maps are continouos. The commmutation then comes done to checking where infinity gets mapped which is always the basepoint of the wedge product.
\end{proof}

\begin{lemma}[Reflection in one variable]\label{lem:reflection-is-minus-one}
	Let $R_1:\R\to\R,\ x\mapsto -x$, and let $Z$ be a pointed compact Hausdorff
	space. Then
	\begin{equation*}
		\left(\overline{R_1}\wedge \id_Z\right)^*=-\id \quad\text{ on } \tilde{K}\left(\overline{\R}\wedge Z\right)\,.
	\end{equation*}
\end{lemma}
\begin{proof}
	Realise $S^1$ as $\R\big/\Z$ with base point $[0]$, so that the map $T$ of
	Corollary \ref{lem:kTheo-inversesInKGivenByHomotopicallyMethods} is
	$T([t])=[-t]$. Define
	\begin{equation*}
		h:\overline{\R}\to \R\big/\Z\,,\quad h(x)=\left[\tfrac{1}{2}+\tfrac{1}{\pi}\arctan(x)\right]\,,\quad h(\infty)=[0]\,.
	\end{equation*}
	Since $x\mapsto \tfrac12+\tfrac1\pi\arctan(x)$ is a strictly increasing bijection
	onto $(0,1)$, the map $h$ is bijective; it is continuous on $\R$ and also at
	$\infty$, because $\arctan(x)\to\pm\tfrac{\pi}{2}$ gives $h(x)\to[1]=[0]$ resp.\
	$h(x)\to[0]$. Being a continuous bijection from a compact space to a Hausdorff
	space, $h$ is a pointed homeomorphism. It intertwines the two reflections:
	\begin{equation*}
		h(-x)=\left[\tfrac12-\tfrac{1}{\pi}\arctan(x)\right]=\left[-\tfrac12-\tfrac{1}{\pi}\arctan(x)\right]=T(h(x))\,,
	\end{equation*}
	using $[-\tfrac12]=[\tfrac12]$, and $h(\overline{R_1}(\infty))=[0]=T(h(\infty))$.
	Hence $h\wedge \id_Z$ is a pointed homeomorphism with
	\begin{equation*}
		(h\wedge\id_Z)\circ\left(\overline{R_1}\wedge\id_Z\right)=(T\wedge 1)\circ (h\wedge \id_Z)\,.
	\end{equation*}
	Applying $\tilde{K}$ and Corollary
	\ref{lem:kTheo-inversesInKGivenByHomotopicallyMethods} gives
	\begin{equation*}
		\left(\overline{R_1}\wedge\id_Z\right)^*=(h\wedge\id_Z)^*\circ (T\wedge1)^*\circ\left((h\wedge\id_Z)^*\right)^{-1}=-\id\,,
	\end{equation*}
	as all maps are ring homomorphisms and thereby commute with multiplicatino by minus one.
\end{proof}




\begin{definition}[K-theoretic orientations]
     We define a \textbf{K-theoretic orientation} of a finite dimensional real vector space $V$ to be a generator of the group $\tilde{K}^{-m}(\overline{V})$. As a sainity check we can confirm: since $\overline{V}\cong S^m$ we are looking at a generator of a group isomorphic to $\tilde{K}(S^{2m})\cong \Z$, which has two possible generators.
\end{definition}








\begin{proposition}\label{prop:linearyTheDeterminandDertermines}
    Let $V$ be an $n$-dimensinoal real vector space and $A\in \GL(V)$. Then 
    \begin{equation*}
        \overline{A}^*=\sign \det(A)\cdot \id \quad \text{ on }\tilde{K}^{-n}(\overline{V})\,.
    \end{equation*}
\end{proposition}
\begin{proof}
    Since the determinant of an endomorphism is invariant under conjugation we can choose a basis and assume that $V=\R^n$. The group $\GL_n(\R)$ as exactly two path components, seperated by the sign of the determinant. Hence, there is a contiuous path $A_{\bullet}:[0,1]\to \GL_n(\R)$  with $A_0=A$ and $A_1=\id$ if $\det(A)>0$ and $R$ if $\det(A)<0$ with $R$ being the reflection $(x_1,\dots,x_m)\mapsto (-x_1,x_2\dots,x_m)$. The map 
    \begin{equation*}
        H:\R^n\times[0,1]\to \R^n\,,\quad H(v,t)\coloneq A_tv
    \end{equation*} is a proper homomorphism into a locally compact space: If $K\subseteq \R^n$ is compact and $A_tv\in K$, then by definition of the operatir norm we have that $\norm{v}\leq (\sup_t\norm{A^{-1}_t})\cdot \max_{k\in K}\norm{k}$. Hence $H^{-1}(K)$ is bounded and closed by continuity. Thereby the the extension $\tilde{H}:\overline{V\times I}\to \overline{V}$ is continuous. Let $q:\overline{V}\times I\to \overline{V\times I}$ be the collaps of $\{\infty\}\times I$ to $\infty$. Then we have that $\overline{H}=\tilde{H}\circ q$ and hence $\overline{H}$ is continuous. If $\det(A)>0$, we are done. Let $\overline{R}$ denote the extension of $R$ to the compactifications. What is left to show is that $\overline{R}^*$ is given by multiplication with minus one.
    For this we write $R=R_1\times \id_{\R^{n-1}}$ where $R_1$ denotes the reflection $x\mapsto -x$ on $\R$. Let $\Xi:\overline{R^n}\to \overline{R}\times \overline{R^{n-1}}$ be the natural homeomorphism from Lemma \ref{lem:orient-compactificationOfProduct}. By the naturality for proper maps we have that
    \begin{equation*}
        \Xi\circ (\overline{R})=(\overline{R_1}\wedge \overline{\id_{\R^{n-1}}})\circ \Xi\,.
    \end{equation*}
    Now $\tilde{K}^{-n}(\overline{R^n})=\tilde{K}(S^n\wedge \overline{R^n})$ and the map $\overline{R}^*$ is induced by the homeomorphism $\id_{S^n}\wedge \overline{R}$. By Lemma \todo{hier lemma, dass wedge kommutative für compakte räume} we have associativity and commutativity in the wedge product letting us define a reordering 
    \begin{equation*} 
        \mathcal{O}:
         S^n\wedge \overline{R^n}
         = S^n\wedge \overline{\R\times \R^{n-1}}
         =S^n\wedge \overline{\R}\wedge \overline{\R^{n-1}} 
         \to 
         \overline{\R}\wedge \underbrace{\big(S^n\wedge \overline{\R^{n-1}}\big)}_{\coloneq Z} 
    \end{equation*} Such that 
    \begin{align*}
        \mathcal{O}\circ (\id_{S^n}\wedge \Xi)(\overline{R})
        &=\mathcal{O}\circ (\id_{S^n}\wedge (\overline{R_1}\wedge \overline{\id_{\R^{n-1}}})\circ \Xi)\\
        &=\mathcal{O}\circ (\id_{S^n}\wedge (\overline{R_1}\wedge \overline{\id_{\R^{n-1}}})) \circ (\id_{S^n}\wedge \Xi) \\
        &=(\overline{R_1}\wedge \id_{S^n}\wedge \overline{\id_{\R^{n-1}}}) \circ \mathcal{O}\circ (\id_{S^n}\wedge \Xi) \\
        &=(\overline{R_1}\wedge \id_Z) \circ \mathcal{O}\circ (\id_{S^n}\wedge \Xi) 
    \end{align*}
    Finally, we define $\Theta\coloneq   \mathcal{O}\circ (\id_{S^n}\wedge \Xi)$ then we can apply Corollary \ref{lem:kTheo-inversesInKGivenByHomotopicallyMethods} to deduce:
    \begin{equation*}
        \overline{R}^*=\Theta^*\circ \underbrace{(\overline{R_1}\wedge \id_Z)}_{=-\id }\circ (\Theta^{-1})^*=-\id\,.
    \end{equation*}
\end{proof}


\begin{comment}

\begin{remark} \label{rem:orient-OrientationsInducesKOrientation}
An orientation of $V$ induces a K-theoretic orientation.
\end{remark}
\begin{proof}
Let $v$ and $v'$ be two equivalent orientations over $V$. Then $M^v_{v'}\simeq \id_V$, since $\GL_{\R}(V)$ has exactly two connected components seperated by the sign of the determinant. (This is a standart theorem for the isomorphic $\GL_m(\R)$). Let $H$ be a homotopy from $M^v_{v'}\simeq \id_V$ given by a continouos path $A:I\to \GL(V)$ with $A(0)=\id$ and $A(1)=M^v_{v'}$.I.e.\ $H:V\times I \to V$ is of the form $H(x,t)=A(t)(x)$. Now extend this homotopy to the following:
\begin{equation*}
	\overline{H}:\overline{V}\times I \to \overline{V}\,,\quad (\infty,t)\mapsto \infty\,.
\end{equation*} Now $H$ is a proper homomorphism into a locally compact space and hence the extension $\tilde{H}:\overline{V\times I}\to \overline{V}$ is continuous. Let $q:\overline{V}\times I\to \overline{V\times I}$ be the collaps of $\{\infty\}\times I$ to $\infty$. Then we have that $\overline{H}=\tilde{H}\circ q$ and hence $\overline{H}$ is continuous. Now let $\kappa_v$ be the coordinate map $V \to \R^m$ that maps each vector to its coordinate representation. Then we have the up to homotopy commutative diagram:
\begin{center}
\begin{tikzcd}
\overline{V} \arrow[r, "\id"] \arrow[d, "\kappa_v"']                                  & \overline{V} \arrow[d, "\kappa_{v'}"] \\
\overline{\R^m} \arrow[r, "M^v_{v'}", bend left] \arrow[r, "\id", bend right=49] & \overline{\R^m}                 
\end{tikzcd}
\end{center}
Hence, the diagram commutes, 
\begin{center}
\begin{tikzcd}
\tilde{K}^{-m}(S^m) \arrow[r, "(\pi_S^{-1})^*"] & \tilde{K}^{-m}(\overline{\R^m}) \arrow[r, "\kappa_v^*"] \arrow[d, "\id"] & \tilde{K}^{-m}(\overline{V}) \arrow[d, "\id"] \\
\tilde{K}^{-m}(S^m) \arrow[r, "(\pi_S^{-1})^*"] & \tilde{K}^{-m}(\overline{\R^m}) \arrow[r, "\kappa^*_{v'}"]               & \tilde{K}^{-m}(\overline{V})                 
\end{tikzcd}
\end{center}
Thereby the induced generators agree.
\end{proof}






\begin{lemma}[K-Theoretic Orientations Detect Orientations]\label{lem:orient-kTheoreticOrientationDetecsOrientation}
    Assume, that $v'$ and $v$ are non-equivalent orientations. Hence, $\det(M^v_{v'})<0$. Then they induce different generators.
\end{lemma}
\begin{proof}
    With a similar argument to above we can conclude, that $M^v_{v'}\simeq \Big(x_1,\dots x_m)\mapsto (x_1,\dots x_{m-1},-x_m)\big)$ and again there extensions to the one point compactifications are homotypic.
    Now we define two maps: 
    \begin{align*}
        T: S^m &\to S^m \\
        (x_1,\dots,x_{m+1})&  \mapsto (-x_1,\dots,x_{m+1})
    \end{align*}
    and
    \begin{align*}
        S:\overline{R^m} & \to \overline{R^m} \\
        (x_1,\dots x_m)  &\mapsto (-x_1,\dots x_m)
    \end{align*}
    With those maps we have the commutative diagram: 
    \begin{center}
        % https://tikzcd.yichuanshen.de/#N4Igdg9gJgpgziAXAbVABwnAlgFyxMJZABgBpiBdUkANwEMAbAVxiRAGUA9AWxAF9S6TLnyEUZAIxVajFmy68BQ7HgJEJ5afWatEIADr6INGACcGWMDGCGASjz79BIDCtHrSU6trl7DxswsrG317bkc+aRgoAHN4IlAAM1MIXkQAJmocCCQyGR02QzQsAH04TmAAWgkI52TUpABmLJzEDXzfA31isorq2qSUtPbs3O9ZXRAAFSdBhsRmkFGM8YK9dn4KPiA
\begin{tikzcd}
S^m                & \overline{\R^m} \arrow[l, "\pi_s^{-1}"]                \\
S^m \arrow[u, "T"] & \overline{\R^m} \arrow[l, "\pi_s^{-1}"] \arrow[u, "S"]
\end{tikzcd}
    \end{center}
    And hence since $S\simeq \overline{M^v_{v'}}$ we get the diagram, is commutative up to homotopy:
\begin{center}
\begin{tikzcd}
S^m                & \overline{\R^m} \arrow[l, "\pi_s^{-1}"]                & \overline{V} \arrow[l, "\kappa_{v'}"]                 \\
S^m \arrow[u, "T"] & \overline{\R^m} \arrow[l, "\pi_s^{-1}"] \arrow[u, "S"] & \overline{V} \arrow[u, "\id"] \arrow[l, "\kappa_{v}"]
\end{tikzcd}
\end{center}
    This induces the commutative diagramm 
\begin{equation} \label{eq: Comutativ diagramm -1 is orientation change}
\begin{tikzcd}
\tilde{K}^{-m}(S^m) \arrow[d, "T^*"'] \arrow[r, "(\pi_S^{-1})^*"] & \tilde{K}^{-m}(\overline{\R^m}) \arrow[d, "S^*"] \arrow[r, "\kappa_{v'}^*"] & \tilde{K}^{-m}(\overline{V}) \arrow[d, "\id"] \\
\tilde{K}^{-m}(S^m) \arrow[r, "(\pi_S^{-1})^*"']                  & \tilde{K}^{-m}(\overline{\R^m}) \arrow[r, "\kappa_v^*"']                & \tilde{K}^{-m}(\overline{V})                 
\end{tikzcd}
    \end{equation}
Now we want to see that $T^*=-\id$.
For this, we start by looking at $K^{-m}(S^m)=K(S^{2m})$. In this case we define 
\begin{align*}
    T': S^m\wedge S^m =S^{2m}       & \to S^{2m}\\
    (x_1,\dots x_m,y_1,\dots, y_m)  & \mapsto (x_1,\dots x_m,-y_1,\dots, y_m)
\end{align*} By definition $T^*=(T')^*$ and we have a homotopy between $T'$ and 
\begin{align*}
    R\wedge \id :S^{2m} & \to       S^{2m}\\
    (x_1,\dots x_{2m})  & \mapsto   (-x_1,\dots x_{2m})
\end{align*} that is given by the rotation $\rho_t$ in the $x_1,y_1$ plane by $t$ degree. This rotation for 90 degree maps $x_1\mapsto y_1$ and $y_1 \mapsto -x_1$, and thereby conjugating $R\wedge \id$ with $\rho$ (for 90 degree) gives $T'$. The homotopy is then given by $\rho_t\circ  R\wedge \id\circ \rho_t^{-1}$. Now the way we defined $S^{2m}\subseteq \R^{2m+1}$ the map $ R\wedge \id$ really is the wedge product of the identity on $S^{2m-1}$ with the map $R$, that exchanges the south and north pole. Hence by the lemma \ref{lem:kTheo-doubleConeFlip} this agrees with the multiplication with $-1$, and hence the map $T^*$ is the multiplication with $-1$. 
Finally, we can compare the induced generators: for this we denote $\beta_v\coloneq (\kappa_v)^*\circ (\pi^{-1}_S)^*(\beta)$ to be the generator given by the oriented basis $v$ and $\beta_{v'}\coloneq (\kappa_{v'})^*\circ (\pi^{-1}_S)^*(\beta)$ given by $v'$. Now we have that 
\begin{align*}
     -\beta_{v'} =\id \circ (\kappa_{v'})^*\circ \underbrace{(\pi^{-1}_S)^*}_{(S^*)^{-1}\circ (\pi^{-1}_S)^*\circ T^*}(-\beta)=\underbrace{\id \circ (\kappa_{v'})^*\circ (S^*)^{-1}}_{=\kappa_v^*}\circ (\pi^{-1}_S)^*\circ \underbrace{T^*(-\beta)}_{=\beta}
    =\beta_v
\end{align*}
\end{proof}
\end{comment}

\begin{corollary}\label{cor:orient-vectorspaceOrientationInducesKtheoreticOrientation}
Any finite dimensional real vector space $V$ is homeomoprhic to $\R^m$ via the coordinate map $\kappa_v$ that maps each vector to its coordinate representation. With this we can introduce a K-theoretic orientation from an ordered basis via the composition: 
   \begin{center}
\begin{tikzcd}
\tilde{K}^{-n}(S^n) \arrow[r, "(\pi_S^{-1})^*"] & \tilde{K}^{-n}(\overline{\R^n}) \arrow[r, "B_v^*"] & \tilde{K}^{-n}(\overline{V})
\end{tikzcd}
   \end{center} With this we can map the bott generator $\beta$ from Definition \ref{def:kTheo-BottGenerator} to a generator of $\tilde{K}^{-n}(\overline{V})$ and hence induce a K-Theoretic orientation. This orientation only depends on the orientation of $V$ and two different orientations induce different K-theoretic orientations
\end{corollary}
\begin{proof}
    The statement follows imidiatly from Proposition \ref{prop:linearyTheDeterminandDertermines}.
\end{proof}
\subsection{Local K-Theoric Orientation}

\begin{lemma}[Minkowsky Gauge]\label{lem:orient-MinkowskyGauge}
    Let $V$ be a finite dimensional real vector space with a norm $\norm{\cdot}$ and let $U\subseteq V$ be open, bounded, convex and contain the origin of $V$. Choose $\epsilon,r>0$ such that
    $\{v\in V\,|\, \norm{v}< \epsilon\}\subseteq U\subseteq \{v\in V\,|\,\norm{v}<r\}$
    and define the Minkowsky gauge
\begin{equation*}
        \mu: V \to [0,\infty]\,,\quad \mu(v)\coloneq \inf\{s>0\,|\, s^{-1}v\in U\}\,.
    \end{equation*} 
    Then:
    \begin{enumerate}[(a)]
        \item $\lambda U \subseteq U$ for $\lambda \in [0,1]$, and $s^{-1}v\in U$ for every $s>\mu(v)$.
        \item $\mu(\lambda v)=\lambda(\mu(v))$ for all $\lambda>0$.
        \item $\frac{\norm{v}}{r}\leq \mu(v)\leq \frac{\norm(v)}{\epsilon}$. In particular $\mu(v)>0$ for $v\neq 0$.
        \item $\mu(v+w)\leq \mu(v)+\mu(w)$.
        \item $\abs{\mu(v)-\mu(w)}\leq \frac{\norm{v-w}}{\epsilon}$, so $\mu$ us Lipschitz and hence continuous;
        \item $U=\{v\in V\,|\,\mu(v)< 1\}$ and $\overline{U}=\{v\in V\,|\,\mu(v)\leq 1 \}$.
    \end{enumerate}
    
\end{lemma}
\begin{proof}
(a) For $x\in U$ and $\lambda\in [0,1]$ we have $\lambda x=\lambda x+(1-\lambda)0\in U$ by convexity and since $0\in U$. Now let $s>\mu(v)$. 
By definition of the infimum there exists a $s_0\in ( \mu(v),s]$ with $s_0^{-1}v\in U$ and $s^{-1}v=\frac{s_0}{s}s_0^{-1}v\in U$ by the first oriperty (since $\frac{s_0}{s}\in [0,1]$).

(b) Follows directly from the definition if one substitutes $s\mapsto \lambda s$.

(c) If $s^{-1}v\in U\subseteq \{v\in V\,|\,\norm{v}\leq r\}$ then $\norm{v}<sr$,i.e.\ $s>\norm{v}r^{-1}$. After the infimum the reads $\mu(v)\geq \norm{v}r^{-1}$. Conversly, if $s>\norm{v}\epsilon^{-1}$ then $\norm{s^{-1}v}<\epsilon$, and hence $s^{-1}v\in U$. Thereby, $\mu(v)\leq \norm{v}\epsilon^{-1}$.

(d) Let $s>\mu(v)$ and $t>\mu(w)$. Then by (a) both $s^{-1}v$ and $t^{-1}w$ lie in $U$, hence by convexity
\begin{equation*}
    \frac{v+w}{s+t}\,=\,\frac{s}{s+t}\cdot \underbrace{\frac{v}{s}}_{\in U}+\underbrace{\frac{w}{t}}_{\in U}\,\in U\,.
\end{equation*} So $\mu(v+w)\leq{s+t}$ and letting $s$ approach $\mu(v)$ and $t$ approach $\mu(w)$ gives the claim.

(e) By (b) and (c), for all $v,w\in V$ we have: $\mu(v)\leq \mu(w)+\mu(v-w)\leq \mu(w)+\norm{v-w}\epsilon^{-1}$. Hence, $\mu(v)-\mu(w)\leq \norm{v-w}\epsilon^{-1}$. Interchangeing $v$ and $w$ gives the symetric statement for $-(\mu(v)-\mu(w))=\mu(w)-\mu(v)\leq \norm{v-w}\epsilon^{-1}$ proofing the statement.

(f) If $v\in U$, then $s^{-1}v \in U$ for some $s<1$, because $U$ is open and $s\mapsto s^{-1}v$ is continuous at $s=1$. Hence $\mu(v)<1$. The other inclusion follows directly from (a): if $\mu(v)<1$ then $1^{-1}(v)\in U$. Thus $U=\{v\in V\,|\, \mu(v) <1\}$. Now, by continuity of $\mu$ the set $\{v\in V|\mu(v)\leq 1\}$ is closed and containes $U$ and thereby also contains $\overline{U}$. conversely if $\mu(v) \leq 1$ then $\mu(tv)=t\mu(v)$ for $t\in(0,1)$ by (b), so $tv\in U$ and if $t$ approaches one, $tv$ approaches $v$. Hence $v$ is in the closure of $U$.
\end{proof}
\begin{lemma}\label{lem:orient-admissableNeighbourhoodsGiveCompactPairs}
    Let $V$ be a finite dimensional real vector space and let $U\subseteq V$ be open, bounded , convex and non empty. Then $\overline{V}\setminus U$ is compact and homeommorphic to the closure $\overline{U}\subseteq V$. In pasrticular $\overline{V}\setminus U $ is contractible and contains $\infty$.
\end{lemma}
\begin{proof}
    The compactnes dolloes from $U$ being open and $\overline{V}$ being comapct. After a translation we may assume that $U$ is open , bounded convex neighbourhood of the origin and hence we can deffine the Minkowsky gauge with resoect to $U$ and after fixing a norm, we may use all properties from Lemma \ref{lem:orient-MinkowskyGauge} (and denote them with just (a)-(f) ). Now we define
    \begin{equation*}
        \iota:\overline{V}\setminus U \to \overline{U}\,,\quad \iota (v)\coloneq \frac{v}{\mu(v)^2}\, (v\neq \infty), \quad \iota(\infty)\coloneq 0. 
    \end{equation*} The well-definition as a map follows from $\mu(v)\geq 1$ if $v\in V\setminus U$ and thereby $\mu(\iota(v))=\mu(v)\mu(v)^2=\mu(v)^{-1}\leq 1$ which concludes the well definition by (f). The map $w\mapsto w\mu(w)^{-2}$ ($0\mapsto \infty$)is a two sided inverse of $\iota$ which can be seen by computation and realizing that $\iota$ maps $\infty$ to zero {lem:orient-MinkowskyGauge} (c). Finally, continuity of $\iota$ in $V\setminus U$ follows from $\mu$ beeing continuous and in $\infty$ the continuity follows from (c):
    \begin{equation*}
        \norm{\iota(v)}=\frac{\norm{v}}{\mu(v)^2} \xrightarrow[\ \norm{v}\to \infty \ ]{} 0\,.
    \end{equation*} The same argument gives us continuouty of our inverse function. Alternitively we can argue that $\iota$ is a continouos bijection from a compact space to a Hausdorr space and thereby a homeomorphism. Finally, $\overline{U}$ is convex and compact and thereby contractible.
\end{proof}

\begin{definition}[Admissible Neighbourhoods]
    Let $V$ be a real finite dimensional vector space. For $x\in V$ let
    \begin{equation*}
        \mathcal{U}(x)\coloneq \{U\subset V \,|\, U \text{ open, bounded, convex neighbourhood of }x \},
    \end{equation*}
    and for $U\in \mathcal{U}(x)$ put $L(U)\coloneq \overline{V}\setminus U$. By Lemma \ref{lem:orient-admissableNeighbourhoodsGiveCompactPairs} $(\overline{V},L(U))$ is a compact pair with $L(U)$ contractible. Furthermore, $(\overline{V},L(U))$ is an NDR pair. To see this, note that in therms of the gauge $\mu$ of $U$ centered at $x$, the space subspace $L(U)$ consists of the points of "radius" greater than or equal to one, together with $\infty$. A radial push towards $L(U)$, damped across a clolor so as to match the identity on $\{v\in V\,|\, \mu(v)\leq 0.5 \}$ and on $\{\infty\}$, gives the required homotopy. The function $u$ from the definition of an NDR pair is adapted to said damping. Notice, how $\mathcal{U}$ is a directed system with respeect to inclusion.
\end{definition}
\begin{definition}[Local K-theoretic orientation] \label{def:orient-localKtheoreticOrientation}
        For $U\in \mathcal{U}(x)$ we call  $K^{-n}(\overline{V},L(U))$ the local group of $V$ at $x$ of size $U$. A \textbf{local K-theoretic orientation of $V$ at $x$} is a generator of this group.
\end{definition}
Next we want to make sure that the naming of said generator is deserved,i.e.\ the generator is induced and detects orientations, and that the size $U$ is irrelevant.

\begin{lemma}[Comparison of Sizes]\label{lem:orient-comparisonOfSizes}
Let $U\subseteq U'$ in $\mathcal{U}(x)$. Then $\id_{\overline{V}}$ gives a map of pairs $(\overline{V},L(U'))\to (\overline{V},L(U))$ and the induced map $\lambda_{U'\,U}:K^{-n}(\overline{V},L(U)) \to K^{-n}(\overline{V},L(U'))$ is an isomorphis. Furthermore let $j$ and $j'$ denote the inclusion of the pair $(\overline{V},\emptyset)$ into $(\overline{V},L(U))$ and $(\overline{V},L(U'))$. Then, the triangle 
\begin{center}
    \begin{tikzcd}
        K^{-n}(\overline{V},L(U)) \arrow[rd,swap,"j^*"]\arrow[rr,"\lambda_{U'\,U}"]&&K^{-n}(\overline{V},L(U')) \arrow[ld,"j'^*"]\\
        &\tilde{K}^{-n}(\overline{V})
    \end{tikzcd}
\end{center} commutes, and all maps are isomorphisms and all groups are isomorphic to $\Z$.
\end{lemma}
\begin{proof}
    The commutation follows from the topological commutation of the maps that induce the diagram. The maps $j$ and $j^*$ induce isomorphisms by Lemma \ref{lem:kTheo-collapsingOfContractibleSubspace}, as they are induced by collapsing contractible spaces where contractability follows from Lemma \ref{lem:orient-admissableNeighbourhoodsGiveCompactPairs}. Hence $\lambda_{U'\, U}$ is also an isomorphism. Finally $K^{-\dim(V)}(\overline{V})\cong \Z$ gets a generator as a K-theoretic orientation.
\end{proof}
\begin{definition}
    For $x\in V$ and $U\in \mathcal{U}(x)$ define the isomorphism
    \begin{equation*}
        \Theta_{x,U}:\tilde{K}^{-n}(\overline{V}) \to K^{-n}(\overline{V},L(U))
    \end{equation*} to be the inverse of the map induced by the pair inclusion $(\overline{V},\emptyset)\to (\overline{V},L(U))$.
\end{definition}
\begin{proposition}[A K-theoretic Orientation Induces a Local Orientation]
    Let $\beta$ be a K-theoretic orientation ov $V$, i.e.\ a generator of $\tilde{K}^{-n}(\overline{V})$. Then for every $x\in V$ and every $U\in \mathcal{U}(x)$ the element $\Theta_{x,U}(\beta)$ is a local K-theoretic orientation of $V$ at $X$ and these are compatible: $\lambda_{U'\,U}\Theta_{x,U}(\beta)=\Theta_{x,U'}(\beta)$ for $U'\subseteq U$.  
\end{proposition}
\begin{proof}
    All statements follow from Lemme \ref{lem:orient-comparisonOfSizes}, as $\Theta_{x,U}$ is an isomorphism and thereby maps generators to generators and the compatibility statement follows from the commutatuve triangle.
\end{proof}
\begin{corollary}
    An orientation of $V$ (i.e.\ an equivalence class of ordered bases) induce via Corollary \ref{cor:orient-vectorspaceOrientationInducesKtheoreticOrientation} a K-theoretic orientation which in turn induces a local K-Theoretic orientation at every $x\in V$.
\end{corollary}

\begin{proposition}[Naturality]\label{prop:orient-localKOrientationIsNatural}
    Let $A:V\to W$ be an isomorphism of $n$-dimenisonal real vector spaces, $x\in V$ and $U\in \mathcal{U}(x)$. Then $A(U)\in \mathcal{U}(Ax)$, the extension $\overline{A}:\overline{V} \to \overline{W}$ is a homeomorphism with $\overline{A}(LU)=L(AU)$, and the square
    \begin{center}
       \begin{tikzcd}
           \tilde{K}^{-n}(\overline{W})  \arrow[r,"\overline{A}^*"]   \arrow[d,"\Theta_{Ax,AU}"] 
         & \tilde{K}^{-n}(\overline{V})  \arrow[d,"\Theta_{x,U}"]                                \\
           K^{-n}(\overline{W},L(AU))    \arrow[r,"\overline{A}^*"]
         & K^{-n}(\overline{V},L(U))
       \end{tikzcd}
    \end{center} commutes. The same holds for affine isomorphisms.
\end{proposition}
\begin{proof}
    A linear isomorphism carries open bounded convex sets to open bounded convex sets, so $A(U)\in \mathcal{U}(AX)$. $A$ is proper, so $\overline{A}$ is a homeomorphism that maps $L(U)$ onto $L(AU)$. If we replace the maps $\Theta$ with the inverses, the corresponding wquare commutes, as its inducing topological square commutes, and thereby our square also commutes.
\end{proof}

\begin{corollary}
    For $A\in \GL(V)$, $x\in V$ and $U\in \mathcal{U}(x)$ the induced map $\overline{A}^*:K^{-n}(\overline{V},L(AU))\to K^{-n}(\overline{V},L(Q))$ equals $\sign \det(A)$ after the canonical identification $\Theta_{x,U}\circ \Theta_{Ax,AU}^{-1}$.
\end{corollary}
\begin{proof}
    This follows directly from Propositions \ref{prop:linearyTheDeterminandDertermines} and \ref{prop:orient-localKOrientationIsNatural}.
\end{proof}
\subsection{Orientations Calculated by a Jacobean Sign}
We now write $B_r \coloneq \{v\in \R^n\,|\, \norm{v}<r\}$ and $D_r\coloneq \overline{B_r}$ for the euclidean balls around $0$, and $L_r\coloneq \overline{R^n}\setminus B_r=L(B_r)$.
\begin{lemma}[Linearisation] \label{lem:orient-Linearisation}
Let $\Omega\subseteq \R^n$ be open with $0\in \Omega$ and let $0\in \Omega$ and let $g:\Omega \to \R^n$ be $C^1$ with 
\begin{equation*}
    g(0)=0\,,\quad g(w)\neq 0 \text{ for }w\in \omega\setminus \{0\}\,,\quad A\coloneq \diff g_0 \in \GL_n(\R)\,.
\end{equation*} Choose $\delta>0$ with $D_\delta\in \Omega$ and let $\rho\in(0,\delta)$. Then there are $c>0$ and $R>c$ such that both $\restr{g}{D_{\delta}}$ and $\restr{A}{D_{\delta}}$ are maps of pairs
\begin{equation*}
    (D_{\delta},D_{\delta}\setminus B_\rho) \to (D_R,D_r\setminus B_c)
\end{equation*} are homotopic as such. In particular $(\restr{g}{D_{\delta}})^*=(\restr{A}{D_{\delta}})^*$ on $K^{-n}(D_R,D_R\setminus B_c)$. Moreover c may be taken arbitrarily small and $R$ arbitrarily large.
\end{lemma}
\begin{proof}
First we define the map 
    \begin{equation*}
        M:D_\delta \to \Mat_n(\R)\,,\quad M(w)\coloneq \int_0^1\restr{\diff g}{s\cdot w}\diff s\,.
    \end{equation*} Then $M$ is continuous, $M(0)=A$ and 
    \begin{equation*}
        g(w)=g(w)-g(0)=\int_0^1 \frac{\diff}{\diff s}(g(s\cdot w))\diff s=M(w)w
    \end{equation*} With this we define the mapping
    \begin{equation*}
        H:D_\delta \times [0,1] \to \R^n\,,\quad H(v,t)\coloneq M(tv)v\,,
    \end{equation*}which is continuous, whith $H(\cdot,1)=\restr{g}{D_\delta}$ and $H(\cdot,0)=\restr{A}{D_\delta}$. For $v\neq 0$ and $t\in (0,1]$ we can calculate
    \begin{equation*}
        H(v,t)=\frac{M(tv)(tv)}{t}=\frac{g(tv)}{t}\neq 0
    \end{equation*} because $g$ only vanishes at the origin. At $t=0$ we have $H(v,0)=Av\neq 0$. Hence, the set $\{H(v,t)\,|\, \rho\leq \norm{v}\leq \delta,\, t\in [0,1]\}$ is compact and does not contain $0$. Thereby the number
    \begin{equation*}
        c_0\coloneq \min\{\norm{H(v,t)}\,|\, \rho\leq \norm{v}\leq \delta,\, t\in [0,1]\}
    \end{equation*} is greater then zero. Now let $c\in (0,c_0]$. Since $H(D_\delta\times [0,1])$ is bounded we may choose $R$ big enough such that $H(D_\delta\times [0,1]) \subseteq B_R$.
    With those constants we can conclude that $H(v,t)\in D_R$ for all $(v,t)$ and if $\norm{v}\geq \rho$ we have $\norm{H(v,t)}\geq c_0\geq c$ meaning that $H(v,t)\in D_r\setminus B_c$. Thereby we have our homotopy of pairs. 
\end{proof}
\begin{proposition}[Local Index] \label{prop:orient-localIndex}
    Let $n\geq 1$ and let $\Psi:\overline{\R^n} \to \overline{\R^n}$ be continuous with
    \begin{enumerate}
        \item $\Psi^{-1}(0)=\{0\}$.
        \item There is an open $\Omega\subseteq \R^n$ with $0\in \Omega$, $\Psi(\Omega)\in \R^n$ and $\restr{\Psi}{\Omega}$ is of class $C^1$.
        \item $A\coloneq \diff(\restr{\Psi}{\omega})_0$ is invertible.
        \end{enumerate} Then 
        \begin{equation*}
            \Psi^*=\sign \det(A) \cdot \id\,,\quad \text{on }\tilde{K}^{-n}(\overline{\R^n})\,.
        \end{equation*}
\end{proposition}
\begin{proof}
    Choose $\delta>0$ such that $D_\delta\subseteq \Omega$ and fix $\rho \in (0,\delta)$. Choose $c,R$ as in Lemma \ref{lem:orient-Linearisation} applied to $\restr{g}{\omega}$ and assume $c$ to be small enough such that 
    \begin{equation}
        \Psi(\overline{R^n}\setminus B_\rho) \cap B_c=\emptyset\,,
    \end{equation} which is possible, as the left side of the intersection doesn't contain the origin and is closed.  Hence, there exist a small ball arround the origin that isn't contained. Furthermore, define $\sigma\coloneq \min\{\norm{Au}\,|\, \norm{u}=1\}>0$. Now assume $c\leq \sigma\cdot \rho$. Then we get two maps of pairs 
    \begin{align*}
        \Psi: &(\overline{R^n},L_\rho)\to (\overline{R^n},L_c)\\
        \overline{A}: &(\overline{R^n},L_\rho)\to (\overline{R^n},L_c)
    \end{align*}For the well definition of the second map assume $\norm{v}\geq \rho$. Then: 
    \begin{equation*}
        \norm{Av}\geq \sigma \norm{v} \geq \sigma \rho \geq c \,.
    \end{equation*}
    Now we apply Lemma \ref{lem:orient-admissableNeighbourhoodsGiveCompactPairs} to get the two isomorphisms induced by inclusions:
    \begin{equation*}
        j_\rho^*:K^{-n}(\overline{\R^n},L_\rho)\to K^{-n}(D_\delta,D_\delta\setminus B_\rho)\,,\quad 
        j_c^*:K^{-n}(\overline{\R^n},L_c)\to K^{-n}(D_\delta,D_\delta\setminus B_c)\,.
    \end{equation*}
    Because $\Psi(D_\delta)\subseteq B_R$ and $\overline{A}(D_\delta)\subseteq B_R$ (by choice of $R$) we get the commuting square
    \begin{center}
        \begin{tikzcd}
            K^{-n}(\overline{R^n},L_c)\arrow[r,"F^*"] \arrow[d,"j_c^*"]&K^{-n}(\overline{R^n},L_\rho)\arrow[d,"j_\rho^*"]\\
            K^{-n}(D_R,D_R\setminus B_c)\arrow[r,"\restr{F}{D_\delta}^*"]&K^{-n}(D_\delta,D_\delta\setminus B_\rho)
        \end{tikzcd}
    \end{center}for both $F=\Psi$ and $F=\overline{A}$. (Commutativity follows from the commutativity of the topological square that induces). Now we can apply Lemma \ref{lem:orient-Linearisation} to deduce that both maps in the bottom row agree. As the vertical maps are isomorphisms we can deduce, that also the top two maps agree, namely $\Psi^*=\overline{A}^*$.
     Finally by lemma \ref{lem:orient-quotientIsCompactification} we now that $L_\rho $ and $L_c$ are contractible and contain $\infty$. Hence the collaps $c_i: \overline{R^n}\to \overline{\R^n}\big/L_i$ for $i=\rho,c$  induces isomorphism that become the vertical maps in the commutative square:
    \begin{center}
         \begin{tikzcd}
            K^{-n}(\overline{R^n},L_c)\arrow[r,"F^*"] \arrow[d,"c_c^*"]&K^{-n}(\overline{R^n},L_\rho)\arrow[d,"c_\rho^*"]\\
            \tilde{K}^{-n}(\overline{R^n})\arrow[r,"F^*"]&\tilde{K}^{-n}(\overline{R^n})
        \end{tikzcd}
    \end{center} Again for $F=\Psi$ and $F=\overline{A}$. Hence, again we can conclude that $\Psi^*=\overline{A}^*$ now as maps in $\tilde{K}^{-n}(\overline{\R^n})$, where Proposition \ref{prop:linearyTheDeterminandDertermines} tells us that $\overline{A}^*=\sign \det(A)\cdot \id$.
\end{proof}
\begin{remark}[Orientation convention]\label{conv:orient-orientationOfTheSphere}
    Let $n\geq 1$. We orient $S^n$ such that the stereographic projection $\pi_S$ from Definition \ref{def:orient-fromRelaSpaceToTheSphere} is orientation preserving.
\end{remark}



\begin{lemma}\label{lem:orient-rotationsPreserveOrientation}
	Every $C\in\SO(n+1)$ acts on $S^n$ as an orientation preserving diffeomorphism,
	and $\pi_S\circ C$ is an orientation preserving chart on $C^{-1}(S^n\setminus\{S\})$.
\end{lemma}
\begin{proof}
	Choose a path $\gamma:[0,1]\to\SO(n+1)$ with $\gamma(0)=\id$ and $\gamma(1)=C$,
	which exists since $\SO(n+1)$ is path-connected. Fix $x\in S^n$. The map
	$t\mapsto \sign\det\left(\restr{\diff \gamma(t)}{x}\right)$ takes values in
	$\{-1,1\}$ and is continuous, because in charts it is the sign of the determinant
	of a matrix depending continuously on $t$, which never vanishes. Hence it is
	constant, and at $t=0$ it equals $+1$. The second claim follows from the chain
	rule together with Corollary \ref{cor:orient-transitionFunctionPositiveDeterminantAlsoOrientationPreserving}.
\end{proof}
\begin{lemma}[Single Regular Preimage]\label{lem:orient-singleRegularPreimage}
Let $n\geq 1$ and let $\xi:S^n\to S^n$ be continuous. Suppose that there are $y,z\in S^n$ with 
\begin{equation*}
    \xi^{-1}(z)=\{x\}\,,
\end{equation*}such that $\xi$ is $C^1$ near $y$ and $\restr{\diff \xi}{y}$ is invertible. Then
\begin{equation*}
    \xi^*=\sign \det(d\xi_y)\cdot \id\,,\quad \text{ on }\tilde{K}^{-n}(S^n)\,.
\end{equation*}
\end{lemma}
\begin{proof}
    Let $S$ be the south pole, $N$ be the north pole and 
    \begin{equation*}
        \overline{\pi_S}:S^n \to \overline{\R^n}
    \end{equation*} be the compactifications of $\pi_S$  to the one-point-compactifications, meaning $\overline{\pi_S}(S)=\infty$. Remember that $\pi_S(N)=0$ is. Now choose two rotations $B,C\in \SO(n+1)$ such that $B(z)=N$ and $C(y)=N$. Then define the two homeomorphisms
    \begin{equation*}
        \overline{\phi}\coloneq \overline{\pi_S}\circ C\,,\quad \overline{ \psi}\coloneq \overline{\pi_S}\circ B
    \end{equation*} that restrict charts over $S^n\setminus C^{-1}(S)$ respectivly over $S^n\setminus B^{-1}(S)$. Those charts are orientation preserving by Lemma \ref{lem:orient-rotationsPreserveOrientation}.
    Now define
    \begin{equation*}
        \Psi:\overline{\psi}\circ \xi \circ \overline{\phi}^{-1}\,:\, \overline{\R^n}\to \overline{\R^n}\,,
    \end{equation*}
    Then $\Psi$ satisfies the hypotheses of Proposition \ref{prop:orient-localIndex}:
\begin{enumerate}
    \item $\Psi(w)=0$ if and only if $\xi\circ \overline{\phi}^{-1}(w)=z$ if and only if $\overline{\phi}^{-1}(w)=y$ which is only the case for $w=0$.
    \item Near $0$ all maps that compose to $\Psi$ are diffeomorphisms and hence $\Psi$ is a diffeomorphism near $0$.
    \item By constructiuon $\restr{\diff \Psi}{0}$ is the coordinate representative of $\restr{\diff \xi}{y}$ in the smooth charts $\phi$ and $\psi$. Invertibility of the latter in the hypothesus exactly states the invertibility of the coordinate representation.
\end{enumerate}
Hence, by Proposition \ref{prop:orient-localIndex} we have that $\psi^*=\sign \det (\restr{\diff \Psi}{0})\cdot \id$ on $\tilde{K}^{-n}(\overline{\R^n})$. Finally, since $\phi$ and $\psi$ are orientation preserving, this sign equials $\sign \det(d\xi_y)$ by Lemma \ref{lem:orient-computingSignInCharts}. Finally, since $SO(n+1)$ is path connected there is a path 
$\gamma:[0,1]\to \SO(n+1)$ with $\gamma(0)=C$ and $\gamma(1)=B$ . Then the map 
    \begin{equation*}
        S^n\times [0,1] \to \overline{\R^n}\,,\quad (x,t)\mapsto \overline{\pi_S}(\gamma(t)x)
    \end{equation*} is a homotopy from $\overline{\phi}$ to $\overline{\psi}$ givin us that $\overline{\phi}^*=\overline{\psi}^*$ and both maps are isomorphisms. Hence:
    \begin{equation*}
        \xi^*=(\overline{\phi}^*)\circ \Psi^*\circ (\overline{\psi}^*)^{-1}=\sign \det(d\xi_y)\circ \underbrace{(\overline{\phi}^*)\circ (\overline{\psi}^*)^{-1}}_{=\id}\,.
    \end{equation*}
\end{proof}